% Options for packages loaded elsewhere
\PassOptionsToPackage{unicode}{hyperref}
\PassOptionsToPackage{hyphens}{url}
\PassOptionsToPackage{dvipsnames,svgnames,x11names}{xcolor}
%
\documentclass[
  letterpaper,
  DIV=11,
  numbers=noendperiod]{scrreprt}

\usepackage{amsmath,amssymb}
\usepackage{iftex}
\ifPDFTeX
  \usepackage[T1]{fontenc}
  \usepackage[utf8]{inputenc}
  \usepackage{textcomp} % provide euro and other symbols
\else % if luatex or xetex
  \usepackage{unicode-math}
  \defaultfontfeatures{Scale=MatchLowercase}
  \defaultfontfeatures[\rmfamily]{Ligatures=TeX,Scale=1}
\fi
\usepackage{lmodern}
\ifPDFTeX\else  
    % xetex/luatex font selection
\fi
% Use upquote if available, for straight quotes in verbatim environments
\IfFileExists{upquote.sty}{\usepackage{upquote}}{}
\IfFileExists{microtype.sty}{% use microtype if available
  \usepackage[]{microtype}
  \UseMicrotypeSet[protrusion]{basicmath} % disable protrusion for tt fonts
}{}
\makeatletter
\@ifundefined{KOMAClassName}{% if non-KOMA class
  \IfFileExists{parskip.sty}{%
    \usepackage{parskip}
  }{% else
    \setlength{\parindent}{0pt}
    \setlength{\parskip}{6pt plus 2pt minus 1pt}}
}{% if KOMA class
  \KOMAoptions{parskip=half}}
\makeatother
\usepackage{xcolor}
\setlength{\emergencystretch}{3em} % prevent overfull lines
\setcounter{secnumdepth}{5}
% Make \paragraph and \subparagraph free-standing
\makeatletter
\ifx\paragraph\undefined\else
  \let\oldparagraph\paragraph
  \renewcommand{\paragraph}{
    \@ifstar
      \xxxParagraphStar
      \xxxParagraphNoStar
  }
  \newcommand{\xxxParagraphStar}[1]{\oldparagraph*{#1}\mbox{}}
  \newcommand{\xxxParagraphNoStar}[1]{\oldparagraph{#1}\mbox{}}
\fi
\ifx\subparagraph\undefined\else
  \let\oldsubparagraph\subparagraph
  \renewcommand{\subparagraph}{
    \@ifstar
      \xxxSubParagraphStar
      \xxxSubParagraphNoStar
  }
  \newcommand{\xxxSubParagraphStar}[1]{\oldsubparagraph*{#1}\mbox{}}
  \newcommand{\xxxSubParagraphNoStar}[1]{\oldsubparagraph{#1}\mbox{}}
\fi
\makeatother


\providecommand{\tightlist}{%
  \setlength{\itemsep}{0pt}\setlength{\parskip}{0pt}}\usepackage{longtable,booktabs,array}
\usepackage{calc} % for calculating minipage widths
% Correct order of tables after \paragraph or \subparagraph
\usepackage{etoolbox}
\makeatletter
\patchcmd\longtable{\par}{\if@noskipsec\mbox{}\fi\par}{}{}
\makeatother
% Allow footnotes in longtable head/foot
\IfFileExists{footnotehyper.sty}{\usepackage{footnotehyper}}{\usepackage{footnote}}
\makesavenoteenv{longtable}
\usepackage{graphicx}
\makeatletter
\newsavebox\pandoc@box
\newcommand*\pandocbounded[1]{% scales image to fit in text height/width
  \sbox\pandoc@box{#1}%
  \Gscale@div\@tempa{\textheight}{\dimexpr\ht\pandoc@box+\dp\pandoc@box\relax}%
  \Gscale@div\@tempb{\linewidth}{\wd\pandoc@box}%
  \ifdim\@tempb\p@<\@tempa\p@\let\@tempa\@tempb\fi% select the smaller of both
  \ifdim\@tempa\p@<\p@\scalebox{\@tempa}{\usebox\pandoc@box}%
  \else\usebox{\pandoc@box}%
  \fi%
}
% Set default figure placement to htbp
\def\fps@figure{htbp}
\makeatother
% definitions for citeproc citations
\NewDocumentCommand\citeproctext{}{}
\NewDocumentCommand\citeproc{mm}{%
  \begingroup\def\citeproctext{#2}\cite{#1}\endgroup}
\makeatletter
 % allow citations to break across lines
 \let\@cite@ofmt\@firstofone
 % avoid brackets around text for \cite:
 \def\@biblabel#1{}
 \def\@cite#1#2{{#1\if@tempswa , #2\fi}}
\makeatother
\newlength{\cslhangindent}
\setlength{\cslhangindent}{1.5em}
\newlength{\csllabelwidth}
\setlength{\csllabelwidth}{3em}
\newenvironment{CSLReferences}[2] % #1 hanging-indent, #2 entry-spacing
 {\begin{list}{}{%
  \setlength{\itemindent}{0pt}
  \setlength{\leftmargin}{0pt}
  \setlength{\parsep}{0pt}
  % turn on hanging indent if param 1 is 1
  \ifodd #1
   \setlength{\leftmargin}{\cslhangindent}
   \setlength{\itemindent}{-1\cslhangindent}
  \fi
  % set entry spacing
  \setlength{\itemsep}{#2\baselineskip}}}
 {\end{list}}
\usepackage{calc}
\newcommand{\CSLBlock}[1]{\hfill\break\parbox[t]{\linewidth}{\strut\ignorespaces#1\strut}}
\newcommand{\CSLLeftMargin}[1]{\parbox[t]{\csllabelwidth}{\strut#1\strut}}
\newcommand{\CSLRightInline}[1]{\parbox[t]{\linewidth - \csllabelwidth}{\strut#1\strut}}
\newcommand{\CSLIndent}[1]{\hspace{\cslhangindent}#1}

\KOMAoption{captions}{tableheading}
\makeatletter
\@ifpackageloaded{bookmark}{}{\usepackage{bookmark}}
\makeatother
\makeatletter
\@ifpackageloaded{caption}{}{\usepackage{caption}}
\AtBeginDocument{%
\ifdefined\contentsname
  \renewcommand*\contentsname{Table of contents}
\else
  \newcommand\contentsname{Table of contents}
\fi
\ifdefined\listfigurename
  \renewcommand*\listfigurename{List of Figures}
\else
  \newcommand\listfigurename{List of Figures}
\fi
\ifdefined\listtablename
  \renewcommand*\listtablename{List of Tables}
\else
  \newcommand\listtablename{List of Tables}
\fi
\ifdefined\figurename
  \renewcommand*\figurename{Figure}
\else
  \newcommand\figurename{Figure}
\fi
\ifdefined\tablename
  \renewcommand*\tablename{Table}
\else
  \newcommand\tablename{Table}
\fi
}
\@ifpackageloaded{float}{}{\usepackage{float}}
\floatstyle{ruled}
\@ifundefined{c@chapter}{\newfloat{codelisting}{h}{lop}}{\newfloat{codelisting}{h}{lop}[chapter]}
\floatname{codelisting}{Listing}
\newcommand*\listoflistings{\listof{codelisting}{List of Listings}}
\makeatother
\makeatletter
\makeatother
\makeatletter
\@ifpackageloaded{caption}{}{\usepackage{caption}}
\@ifpackageloaded{subcaption}{}{\usepackage{subcaption}}
\makeatother

\ifLuaTeX
\usepackage[bidi=basic]{babel}
\else
\usepackage[bidi=default]{babel}
\fi
\babelprovide[main,import]{arabic}
% get rid of language-specific shorthands (see #6817):
\let\LanguageShortHands\languageshorthands
\def\languageshorthands#1{}
\ifPDFTeX
  \TeXXeTstate=1
  \newcommand{\RL}[1]{\beginR #1\endR}
  \newcommand{\LR}[1]{\beginL #1\endL}
  \newenvironment{RTL}{\beginR}{\endR}
  \newenvironment{LTR}{\beginL}{\endL}
\fi
\usepackage{bookmark}

\IfFileExists{xurl.sty}{\usepackage{xurl}}{} % add URL line breaks if available
\urlstyle{same} % disable monospaced font for URLs
\hypersetup{
  pdftitle={Retraites},
  pdfauthor={Mahdi Ben Jelloul},
  pdflang={ar},
  colorlinks=true,
  linkcolor={blue},
  filecolor={Maroon},
  citecolor={Blue},
  urlcolor={Blue},
  pdfcreator={LaTeX via pandoc}}


\title{Retraites}
\author{Mahdi Ben Jelloul}
\date{}

\begin{document}
\maketitle

\renewcommand*\contentsname{Table of contents}
{
\hypersetup{linkcolor=}
\setcounter{tocdepth}{2}
\tableofcontents
}

\bookmarksetup{startatroot}

\chapter{Le système de retraite
tunisien}\label{le-systuxe8me-de-retraite-tunisien}

\section{Présentation générale}\label{pruxe9sentation-guxe9nuxe9rale}

Il existe deux caisses nationales gérant les régimes dispensant une
retraite à leurs affiliés.

La Caisse nationale de retraite et de prévoyance sociale (CNRPS) gère
les régimes de la fonction publique civile et militaire, des régimes
spéciaux des députés, des gouverneurs et celui des membres du
gouvernement.

La Caisse nationale de sécurité sociale (CNSS) gère les régimes assurant
une etraite aux travailleurs du secteur privé.

\bookmarksetup{startatroot}

\chapter{Le secteur public}\label{le-secteur-public}

\subsection{Avant l'indépendance}\label{avant-linduxe9pendance}

Le premier régime de sécurité sociale a été instauré en Tunisie en 1898
avec la création de la Société de prévoyance des fonctionnaires et
employés tunisiens (SPFET).

En 1948, un régime de retraite pour les personnels des compagnies
d'électricité, de gaz et des transports a été crée (Décret du 26 août
1948).

\subsection{Les pensions civiles et
militaires}\label{les-pensions-civiles-et-militaires}

Depuis 1985 le texte de référence est la Loi 85-12 du 5 Mars 1985
portant régime des pensions civiles et militaires de retraite et des
survivants dans le secteur public\footnote{L'ouvrage de (Sdiri 2014) est
  une bonne référence pour le fonctionnement du sstème de retraite en
  vigueur en 2012 soit avec l'essentiel des disposition de la loi de
  1985.}.

Citation de loi (\emph{Loi n° 2016-1321}, د.ت)

Elle couvre tous les agents du secteur public quels que soient leur
situation administrative, les modalités de paiement de leur
rémunération, leux sexe, leur nationalité. Sont concernés, les agents de
l'État, des collectivités publiques locales, les établissements publics
à caractère administratif, les établissements publics à caractère
industriel et commercial et les sociétés nationales.

\subsubsection{Les conditions du droit à
pension}\label{les-conditions-du-droit-uxe0-pension}

\paragraph{Conditions d'âge}\label{conditions-duxe2ge}

L'âge légal de départ à la retraite a été fixé à 60 ans.

Il est abaissé à 55 ans pour les agents des cadres actifs ainsi que pou
les ouvriers accomplissant des tâches pénibles, dans des conditions
insalbubres ou astreignantes.

Il peut être prorogé jusqu'à 65 ans par décret.

La loi n° 2009-20 du 13 avril 2009, portant dispositions exceptionnelles
relatives à la retraite de certains professeurs de l'enseignement
supérieur porte l'âge de mise à la retraite à 65 ans. Les enseignants
concernés sont:

\begin{itemize}
\item
  les professeurs de l'enseignement supérieur et les maîtres de
  conférences de l'enseignement supérieur rattachés aux établissements
  universitaires et aux établissements de recherche scientifique civils
  et militaires
\item
  les professeurs hospitalo-universitaires et les maîtres de conférences
  agrégés hospitalo-universitaires
\end{itemize}

Il peut être prorogé jusqu'à 70 ans par décret.

L'âge légal été modifé par la loi du 30 avril 2019 pour être porté à 62
ans. Cette augmentation est graduelle. Elle est portée à 61 au 1er
juiller 2019 et à 62 ans au 1er janvier 2020.

\paragraph{Condition de stage}\label{condition-de-stage}

La durée de stage minimale qui doit être effectuée pour l'ouverture des
droits est de 15 ans de service.

Cette durée est réduite à 10 ans pour les ouvriers.

\subsubsection{Le taux de liquidation}\label{le-taux-de-liquidation}

\subsubsection{Le salaire de
référence}\label{le-salaire-de-ruxe9fuxe9rence}

L'assiette de cotisation comprend l'ensemble des éléments permanents de
la rémunération de l'agent.

Le ssalaire de référence correspond à :

\begin{itemize}
\tightlist
\item
  la dernière rémunération ayant l'objet de retenues pour pension
  pendant au moins trois ans
\item
  ou la rémunération afférente à la fonction la plus élevée
  effectivement exercée pendant au moins deux ans au cours de la
  carrière de l'agent.
\end{itemize}

\subsubsection{Les départs anticipés}\label{les-duxe9parts-anticipuxe9s}

Le droit à pension de retraite peut s'acquérir avant d'atteindre l'âge
légal dans les cas suivants :

\begin{itemize}
\item
  En cas d'invalidité
\item
  Sur demande de l'assuré et après accord de l'employeur soumis au visa
  du Premier ministre (vu l'accroissement des demandes et des
  répercussions financières)
\item
  En cas de démission
\item
  En cas de révocation par l'employeur pour insuffisance professionelle
\item
  Sur demande des mères ayant 3 enfants de moins de 20 ans ou comprenant
  un enfant ayant un handicap lourd (sans limite d'âge) après accord du
  Premier ministre
\item
  En cas de mise à la retraite d'office après 15 nas de service civils
  et militaires effectifs
\item
  Par arrêté de premier ministre après avis de commission
  d'assainissement et de restructuration des entreprises à participation
  publiques (CAREP)
\end{itemize}

Cependant la date jouissance de la pension ne correspond pas forcément à
la date de départ à la retraite. La joussiance de la pension
n'intervient

\subsubsection{La pension minimale}\label{la-pension-minimale}

La pension minimale s'établit à 2/3 du Smig.

\subsubsection{L'allocation vieillesse}\label{lallocation-vieillesse}

Les agents ayant une durée d'assurance supérieure ou égale à 5 ans et
strictement inférieure à la durée de stage requise par leur statut
peuvent opter pour le remboursement de leurs cotisations retraites ou
demander à jouir d'une allocation vieillesse.

L'allocation vieillesse 50\% du Smig du régime de 2400 heures par an.

\subsubsection{Droits annexes}\label{droits-annexes}

\subsubsection{Rente d'invalidité}\label{rente-dinvalidituxe9}

La pension maximale servie comprise comme la somme de la rente
d'invalidité et de la pension de retraite ne saurait être à 100 \% du
traitement ayant servi de référence pour la liquidation de la pension.

\subsubsection{Les militaires}\label{les-militaires}

\begin{verbatim}
   


`<!-- quarto-file-metadata: eyJyZXNvdXJjZURpciI6Ii4ifQ== -->`{=html}

```{=html}
<!-- quarto-file-metadata: eyJyZXNvdXJjZURpciI6Ii4iLCJib29rSXRlbVR5cGUiOiJjaGFwdGVyIiwiYm9va0l0ZW1OdW1iZXIiOjMsImJvb2tJdGVtRmlsZSI6Il9zZWN0ZXVyX3ByaXZlLnFtZCIsImJvb2tJdGVtRGVwdGgiOjB9 -->
\end{verbatim}

\bookmarksetup{startatroot}

\chapter{Le secteur privé}\label{le-secteur-privuxe9}

\subsection{Le régime des salariés
non-agricoles}\label{le-ruxe9gime-des-salariuxe9s-non-agricoles}

\subsection{Le régime des salariés
agricoles}\label{le-ruxe9gime-des-salariuxe9s-agricoles}

\subsection{Le régime des salariés agricoles
amélioré}\label{le-ruxe9gime-des-salariuxe9s-agricoles-amuxe9lioruxe9}

\subsection{Le régime des travailleurs
non-salariés}\label{le-ruxe9gime-des-travailleurs-non-salariuxe9s}

\subsection{Le régime des travailleurs tunisiens à
l'étranger}\label{le-ruxe9gime-des-travailleurs-tunisiens-uxe0-luxe9tranger}

\subsection{Le régime des travailleurs à bas
revenu}\label{le-ruxe9gime-des-travailleurs-uxe0-bas-revenu}

\subsection{Le régime des artistes, des créateurs et des
intellectuels}\label{le-ruxe9gime-des-artistes-des-cruxe9ateurs-et-des-intellectuels}

\subsection{Le régime complémentaire de
retraite}\label{le-ruxe9gime-compluxe9mentaire-de-retraite}

\subsection{Le régime des étudiants et des
staiaires}\label{le-ruxe9gime-des-uxe9tudiants-et-des-staiaires}

\phantomsection\label{refs}
\begin{CSLReferences}{1}{0}
\bibitem[\citeproctext]{ref-loi2016}
\emph{Loi n° 2016-1321}. د.ت.
\url{https://www.legifrance.gouv.fr/jorf/id/JORFTEXT000033202746}.

\bibitem[\citeproctext]{ref-sdiri2014}
Sdiri, Khaled. 2014. \emph{Analyse économique et pilotage des régimes de
retraite}.

\end{CSLReferences}




\end{document}
